% =============================================================================
% Manual P04 – Anemometría Láser Doppler (LDA/LDV)
% Técnicas de Medida – Ingeniería Aeronáutica
% =============================================================================
\documentclass[12pt,a4paper]{article}

\usepackage[utf8]{inputenc}
\usepackage[T1]{fontenc}
\usepackage[spanish]{babel}
\usepackage[top=2.5cm,bottom=2.5cm,left=3cm,right=2.5cm]{geometry}
\usepackage{amsmath,amssymb,amsfonts}
\usepackage{siunitx}
\sisetup{output-decimal-marker={,}, group-separator={.}}
\usepackage{graphicx,xcolor,float,caption,subcaption}
\usepackage{tikz}
\usetikzlibrary{arrows.meta,positioning,calc,shapes.geometric,
                decorations.markings,patterns,fit,backgrounds,
                decorations.pathmorphing,optics}
\usepackage{booktabs,multirow,array}
\usepackage{listings}
\lstset{
  basicstyle=\ttfamily\small,
  keywordstyle=\color{blue!70!black}\bfseries,
  commentstyle=\color{green!50!black},
  stringstyle=\color{red!60!black},
  backgroundcolor=\color{gray!10},
  frame=single, breaklines=true,
  numbers=left, numberstyle=\tiny\color{gray}
}
\usepackage[colorlinks=true,linkcolor=blue!70!black,citecolor=green!50!black,
            urlcolor=blue!50!black]{hyperref}
\usepackage[backend=biber,style=ieee,sorting=none]{biblatex}
\addbibresource{references_P04.bib}
\usepackage{fancyhdr}
\pagestyle{fancy}
\fancyhf{}
\fancyhead[L]{\small Técnicas de Medida -- Ingeniería Aeronáutica}
\fancyhead[R]{\small P04: LDA -- Perfil de Velocidades}
\fancyfoot[C]{\thepage}

\definecolor{aeroblue}{RGB}{0,51,102}
\definecolor{aerored}{RGB}{153,0,0}
\definecolor{aerogreen}{RGB}{0,102,51}
\definecolor{aerolaser}{RGB}{0,180,0}

% =============================================================================
\begin{document}

\begin{titlepage}
  \centering
  \vspace*{1cm}
  {\color{aeroblue}\rule{\linewidth}{2pt}}\\[0.4cm]
  {\Large\bfseries\color{aeroblue} TÉCNICAS DE MEDIDA EN INGENIERÍA AERONÁUTICA}\\[0.3cm]
  {\color{aeroblue}\rule{\linewidth}{1pt}}\\[1.5cm]
  {\huge\bfseries P04 -- Anemometría Láser\\[0.3cm]Doppler (LDA/LDV)}\\[0.5cm]
  {\Large Perfil de Velocidades en Túnel de Viento}\\[0.8cm]
  {\large Manual de Laboratorio}\\[0.5cm]
  {\large\itshape Ingeniería Aeronáutica}\\[1.5cm]
  \begin{tikzpicture}[scale=0.9,>=Stealth]
    % Two laser beams crossing
    \draw[aerolaser,ultra thick] (0,1.0) -- (3.0,0);
    \draw[aerolaser,ultra thick] (0,-1.0) -- (3.0,0);
    % Measurement volume (ellipse)
    \fill[aerolaser!30] (3.0,0) ellipse (0.25 and 0.12);
    \draw[aerolaser,thick] (3.0,0) ellipse (0.25 and 0.12);
    % Fringes inside volume
    \foreach \dx in {-0.2,-0.1,0.0,0.1,0.2} {
      \draw[aerolaser!60,thin] (3.0+\dx,-0.12) -- (3.0+\dx,0.12);
    }
    % Particle
    \fill[gray!80] (3.0,0) circle (0.06);
    \draw[->,blue!60,thick] (2.5,0) -- (3.5,0)
          node[above,font=\small] {Partícula};
    % Scattered light to detector
    \draw[gray,dashed,thick,->] (3.0,0.12) .. controls (3.5,1.5) .. (5.0,1.5)
          node[right,font=\small] {$f_D$};
    % Beam labels
    \node[aerolaser,above,font=\small] at (1.0,0.7) {Haz 1};
    \node[aerolaser,below,font=\small] at (1.0,-0.7) {Haz 2};
    \node[aerolaser,font=\small] at (3.0,-0.35) {Vol. medida};
    % Fringe spacing
    \draw[<->,aerolaser!60,thin] (2.75,-0.3) -- (3.25,-0.3)
          node[midway,below,font=\tiny] {$\delta_f$};
    % Laser source
    \draw[fill=aeroblue!20,draw=aeroblue,thick] (-1.5,-0.2) rectangle (0,0.2);
    \node[font=\tiny\bfseries] at (-0.75,0) {Láser};
    \draw[->,aerolaser,thick] (0,0.1) -- (0.2,0.8);
    \draw[->,aerolaser,thick] (0,-0.1) -- (0.2,-0.8);
  \end{tikzpicture}\\[1.5cm]
  {\large\textbf{Departamento de Ingeniería Aeroespacial}}\\[0.2cm]
  {\normalsize Laboratorio de Técnicas de Medida}\\[0.5cm]
  {\normalsize \today}
  \vfill
  {\color{aeroblue}\rule{\linewidth}{2pt}}
\end{titlepage}

\tableofcontents
\newpage

% =============================================================================
\section{Introducción}
\label{sec:intro}

La \textbf{Anemometría Láser Doppler} (LDA, también denominada LDV --
Laser Doppler Velocimetry) es una técnica óptica no intrusiva para la medición
puntual de velocidad en flujos de gas y líquido \cite{durst1981}.  A diferencia
del hilo caliente, la LDA no perturba el flujo y puede medir en regiones de
alta temperatura, flujos multifásicos o en presencia de superficies sólidas
próximas.

\subsection{Aplicaciones en ingeniería aeronáutica}

\begin{itemize}
  \item \textbf{Perfiles de velocidades} en secciones de prueba de túneles de
        viento: capas límite, estelas, flujos en difusores.
  \item \textbf{Medición en torno a perfiles alares}: detalle de capa límite,
        burbujas de separación, puntos de reenganche.
  \item \textbf{Aerodinámica interna de motores}: flujo en compresor axial,
        toberas, cámaras de combustión con ventanas ópticas.
  \item \textbf{Rotores de helicóptero}: medición en el plano del rotor con
        seeding de humo.
  \item \textbf{Validación de CFD}: perfiles de $U(y)$ y $V(y)$ como dato de
        referencia para códigos de simulación.
\end{itemize}

% =============================================================================
\section{Objetivos}
\label{sec:objetivos}

\begin{enumerate}
  \item Comprender el principio del efecto Doppler aplicado a la dispersión de
        luz por partículas trazadoras.
  \item Calcular la separación de franjas $\delta_f$ del volumen de medida y
        la velocidad a partir de la frecuencia Doppler $f_D$.
  \item Procesar los archivos de datos LDA del laboratorio (\texttt{files/P4/})
        para obtener el perfil de velocidades $U(y)$.
  \item Calcular estadísticos de turbulencia por posición: $\bar{U}$,
        $\sigma_U$, intensidad de turbulencia $I_u$.
  \item Comparar el perfil medido con las leyes teóricas de capa límite
        (parabólica o logarítmica).
\end{enumerate}

% =============================================================================
\section{Fundamento Teórico}
\label{sec:teoria}

\subsection{Efecto Doppler}

Cuando una fuente de ondas se mueve relativa a un observador, la frecuencia
percibida difiere de la emitida.  En LDA, la \emph{partícula trazadora}
dispersa luz de dos haces que inciden sobre ella con ángulos opuestos $\pm\theta/2$.
La frecuencia Doppler diferencial resulta:

\begin{equation}
  f_D = \frac{2\sin(\theta/2)}{\lambda}\,U_\perp
  \label{eq:fD_general}
\end{equation}

donde $\lambda$ es la longitud de onda del láser y $U_\perp$ la componente de
velocidad perpendicular a la bisectriz de los haces.

\subsection{Modelo de franjas (Fringe Model)}

Los dos haces coherentes se cruzan formando un volumen de medida donde
interfieren constructivamente y destructivamente, creando un patrón de franjas
paralelas (franjas de interferencia).  La separación entre franjas:

\begin{equation}
  \boxed{\delta_f = \frac{\lambda}{2\sin(\theta/2)}}
  \label{eq:fringe_spacing}
\end{equation}

La velocidad de la partícula que atraviesa las franjas:
\begin{equation}
  \boxed{U = f_D \cdot \delta_f = f_D\,\frac{\lambda}{2\sin(\theta/2)}}
  \label{eq:U_LDA}
\end{equation}

La señal de velocidad es independiente de la longitud de la partícula y de la
intensidad de la luz dispersada; solo depende de $f_D$ y de la geometría óptica.

\subsection{Volumen de medida}

Los dos haces gaussianos de diámetro $d_b$ y longitud de onda $\lambda$
enfocados con una lente de focal $f_L$ producen un volumen de medida elipsoidal:

\begin{align}
  d_{mv} &= \frac{4 f_L \lambda}{\pi d_b}
           \quad\text{(diámetro del volumen)} \label{eq:dmv}\\
  l_{mv} &= \frac{d_{mv}}{\sin(\theta/2)}
           \quad\text{(longitud del volumen)} \label{eq:lmv}
\end{align}

El número de franjas $N_f$ en el volumen de medida:
\begin{equation}
  N_f = \frac{d_{mv}}{\delta_f}
  \label{eq:Nf}
\end{equation}

\subsection{Ejemplo numérico: sistema LDA del laboratorio}

\begin{table}[H]
  \centering
  \caption{Parámetros ópticos del sistema LDA del laboratorio.}
  \label{tab:params_LDA}
  \begin{tabular}{lll}
    \toprule
    \textbf{Parámetro} & \textbf{Símbolo} & \textbf{Valor} \\
    \midrule
    Longitud de onda (Ar-Ion)   & $\lambda$     & \SI{514.5}{\nm} \\
    Ángulo de cruce semi-ángulo & $\theta/2$    & \SI{4.5}{\degree} \\
    Focal de la lente           & $f_L$         & \SI{250}{\mm} \\
    Diámetro del haz            & $d_b$         & \SI{1.2}{\mm} \\
    \bottomrule
  \end{tabular}
\end{table}

\paragraph{Separación de franjas:}
\begin{align}
  \delta_f &= \frac{\lambda}{2\sin(\theta/2)}
            = \frac{514.5\times10^{-9}}{2\sin(4.5°)}
            = \frac{514.5\times10^{-9}}{2\times0.07846} \notag\\
           &= \frac{514.5\times10^{-9}}{0.15692}
            \approx \SI{3.279}{\micro\metre}
  \label{eq:delta_f_num}
\end{align}

\paragraph{Velocidad a partir de $f_D$:}
Para una señal Doppler con $f_D = \SI{1.53}{\MHz}$:
\begin{equation}
  U = f_D \cdot \delta_f
    = 1.53\times10^6 \times 3.279\times10^{-6}
    \approx \SI{5.017}{\metre\per\second}
  \label{eq:U_num}
\end{equation}

\paragraph{Volumen de medida:}
\begin{align}
  d_{mv} &= \frac{4\times250\times10^{-3}\times514.5\times10^{-9}}
                 {\pi\times1.2\times10^{-3}}
          = \frac{5.145\times10^{-7}}{3.770\times10^{-3}}
          \approx \SI{136.5}{\micro\metre} \label{eq:dmv_num}\\
  l_{mv} &= \frac{136.5}{\sin(4.5°)} = \frac{136.5}{0.0785}
          \approx \SI{1.74}{\mm} \label{eq:lmv_num}\\
  N_f    &= \frac{136.5}{3.279} \approx 42 \text{ franjas}
  \label{eq:Nf_num}
\end{align}

\subsection{Configuración óptica del sistema LDA}

\begin{figure}[H]
  \centering
  \begin{tikzpicture}[>=Stealth, scale=1.05]
    % Laser source
    \draw[fill=aeroblue!20,draw=aeroblue,very thick] (-1,-0.4) rectangle (1,0.4);
    \node[font=\small\bfseries] at (0,0) {Ar-Ion\\514.5 nm};
    % Beam splitter
    \draw[fill=gray!30,draw=black,thick] (2,-0.4) rectangle (2.2,0.4);
    \node[below,font=\tiny] at (2.1,-0.5) {Divisor};
    \draw[aerolaser,very thick,->] (1,0) -- (2,0);
    % Two beams
    \draw[aerolaser,very thick] (2.2,0.1) -- (4.5,0.9);
    \draw[aerolaser,very thick] (2.2,-0.1) -- (4.5,-0.9);
    % Focusing lens
    \draw[fill=cyan!20,draw=black,thick]
      (4.5,-1.2) -- (4.7,0) -- (4.5,1.2) -- (4.3,1.2) -- (4.5,0) -- (4.3,-1.2) -- cycle;
    \node[below,font=\tiny] at (4.5,-1.3) {Lente $f_L=\SI{250}{\mm}$};
    % Measurement volume
    \fill[aerolaser!30] (6.0,0) ellipse (0.15 and 0.08);
    \draw[aerolaser,thick] (6.0,0) ellipse (0.15 and 0.08);
    % Fringes
    \foreach \dx in {-0.1,-0.05,0.0,0.05,0.1} {
      \draw[aerolaser!50,very thin] (6+\dx,-0.08) -- (6+\dx,0.08);
    }
    % Particle path
    \draw[->,blue!60,thick] (5.2,-0.4) -- (6.8,-0.4)
          node[right,font=\small] {Partícula};
    \draw[blue!60,dashed] (6.0,-0.4) -- (6.0,0);
    % Scattered light
    \draw[gray,dashed,->] (6.0,0.08) .. controls (6.5,1.0) .. (8.0,1.0);
    % Receiving optics
    \draw[fill=cyan!20,draw=black,thick]
      (8.0,0.6) -- (8.2,1.0) -- (8.0,1.4) -- (7.8,1.4) -- (8.0,1.0) -- (7.8,0.6) -- cycle;
    \node[above,font=\tiny] at (8.0,1.5) {Lente rec.};
    % Photomultiplier
    \draw[fill=gray!30,draw=black] (8.5,0.7) rectangle (9.5,1.3);
    \node[font=\tiny] at (9.0,1.0) {PMT};
    \draw[->,gray!50] (8.0,1.0) -- (8.5,1.0);
    % Signal processor
    \draw[fill=aeroblue!20,draw=aeroblue] (9.8,0.6) rectangle (11.0,1.4);
    \node[font=\tiny,aeroblue] at (10.4,1.0) {BSA\\$f_D$};
    \draw[->] (9.5,1.0) -- (9.8,1.0);
    \draw[->,aeroblue,thick] (11.0,1.0) -- (12.0,1.0)
          node[right,font=\small] {$U = f_D\,\delta_f$};
    % Labels
    \node[above,font=\small,aerolaser] at (5.3,0.8) {Haz 1};
    \node[below,font=\small,aerolaser] at (5.3,-0.8) {Haz 2};
    \node[below,font=\small,aerolaser] at (6.0,-0.15) {Vol. medida};
    % Beam crossing angle
    \draw[gray,thin] (4.5,0) -- (6.0,0);
    \pgfmathsetmacro{\ang}{atan(0.9/2.3)}
    \draw[gray,thin] (5.5,0) arc (0:\ang:0.5);
    \node[gray,font=\tiny] at (5.6,0.12) {$\theta/2$};
  \end{tikzpicture}
  \caption{Configuración óptica del sistema LDA: el láser Ar-Ion se divide
           en dos haces que se enfocan con la lente de transmisión sobre el
           volumen de medida.  La luz dispersada se recoge con la óptica
           receptora y se procesa en el Burst Spectrum Analyser (BSA) para
           obtener $f_D$ y, por tanto, $U$.}
  \label{fig:LDA_optics}
\end{figure}

\subsection{Desplazamiento de frecuencia (Bragg cell)}

Para distinguir flujos positivos de negativos, se introduce un desplazamiento
de frecuencia $f_s$ (típicamente \SI{40}{\MHz}) en uno de los haces mediante
una \emph{celda de Bragg}.  La frecuencia efectiva observada:
\begin{equation}
  f_\mathrm{obs} = f_s + f_D
  \label{eq:bragg}
\end{equation}

Si $U > 0$: $f_\mathrm{obs} > f_s$; si $U < 0$: $f_\mathrm{obs} < f_s$.
Esto elimina la ambigüedad de dirección.

% =============================================================================
\section{Estructura de Datos del Laboratorio}
\label{sec:datos}

Los archivos de datos se encuentran en \texttt{files/P4/} y están organizados
en 4 carpetas (FX01G00 a FX04G00), cada una correspondiente a una posición $y$
diferente en el perfil vertical del túnel.  Cada carpeta contiene $\approx 15$
archivos con $\approx 2000$ mediciones cada uno.

\begin{table}[H]
  \centering
  \caption{Estructura de archivos LDA del laboratorio.}
  \label{tab:archivos}
  \begin{tabular}{llll}
    \toprule
    \textbf{Carpeta} & \textbf{Posición y} & \textbf{Archivos} &
    \textbf{Mediciones} \\
    \midrule
    FX01G00 & $y_1$ (inferior) & 15 ficheros & $\sim$30\,000 \\
    FX02G00 & $y_2$ & 15 ficheros & $\sim$30\,000 \\
    FX03G00 & $y_3$ & 15 ficheros & $\sim$30\,000 \\
    FX04G00 & $y_4$ (superior) & 15 ficheros & $\sim$30\,000 \\
    \bottomrule
  \end{tabular}
\end{table}

\textbf{Formato de los archivos} (5 columnas, sin encabezado):
\begin{enumerate}
  \item Índice de muestra
  \item \textbf{Velocidad $U$ [\si{\metre\per\second}]}: componente principal del flujo
  \item Dato auxiliar (tiempo o frecuencia)
  \item \textbf{Velocidad $V$ [\si{\metre\per\second}]}: componente transversal
  \item Indicador de calidad/validación
\end{enumerate}

% =============================================================================
\section{Procesamiento de Datos}
\label{sec:procesamiento}

\subsection{Filtrado de outliers (método IQR)}

Para cada posición $y_k$, se eliminan las mediciones fuera del rango:
\begin{equation}
  [Q_1 - 1.5\,\mathrm{IQR},\; Q_3 + 1.5\,\mathrm{IQR}]
  \label{eq:IQR}
\end{equation}

con $\mathrm{IQR} = Q_3 - Q_1$ (rango intercuartílico) y $Q_1$, $Q_3$ los
percentiles 25 y 75 de la muestra.

\subsection{Estadísticos por posición}

\begin{align}
  \bar{U}(y_k) &= \frac{1}{N_k}\sum_{i=1}^{N_k} U_i
  \label{eq:Umean_y}\\
  \sigma_U(y_k) &= \sqrt{\frac{1}{N_k-1}\sum_{i=1}^{N_k}(U_i-\bar{U})^2}
  \label{eq:sigma_y}\\
  I_u(y_k) &= \frac{\sigma_U(y_k)}{\bar{U}(y_k)}\times100\%
  \label{eq:Iu_y}
\end{align}

\subsection{Ejemplo numérico: estadísticos LDA}

\begin{table}[H]
  \centering
  \caption{Estadísticos LDA por posición (valores representativos).}
  \label{tab:estadisticos}
  \begin{tabular}{lSSSS}
    \toprule
    \textbf{Carpeta} & {$y$ (\si{\mm})} & {$\bar{U}$ (\si{\m\per\s})} &
    {$\sigma_U$ (\si{\m\per\s})} & {$I_u$ (\si{\percent})} \\
    \midrule
    FX01G00 &  25 & 18.3 & 1.83 & 10.0 \\
    FX02G00 &  75 & 22.1 & 1.77 &  8.0 \\
    FX03G00 & 150 & 25.8 & 1.55 &  6.0 \\
    FX04G00 & 225 & 24.6 & 1.72 &  7.0 \\
    \bottomrule
  \end{tabular}
\end{table}

\subsection{Comparación con perfil teórico}

Para flujo turbulento entre placas (canal), el perfil parabólico de Poiseuille
(válido en régimen laminar) es:
\begin{equation}
  U(y) = U_\mathrm{max}\left[1 - \left(\frac{y - h/2}{h/2}\right)^2\right]
  \label{eq:Poiseuille}
\end{equation}

Para flujo turbulento desarrollado, el perfil de la $1/7$ potencia:
\begin{equation}
  \frac{U(y)}{U_\mathrm{max}} = \left(\frac{y}{\delta}\right)^{1/7}
  \label{eq:1_7_power}
\end{equation}

El ajuste del perfil LDA al modelo teórico se evalúa con el coeficiente $R^2$
(ec.~\eqref{eq:R2} de P02).

\subsection{Incertidumbre de la medida LDA}

La incertidumbre de la velocidad medida por LDA proviene principalmente de la
incertidumbre en la frecuencia Doppler $u_{f_D}$:
\begin{equation}
  u_U = \delta_f \cdot u_{f_D}
  \label{eq:uU_LDA}
\end{equation}

Para un Burst Spectrum Analyser (BSA) con resolución $\Delta f = \SI{10}{\kHz}$:
\begin{equation}
  u_U = 3.279\times10^{-6} \times 10\times10^3
      \approx \SI{0.033}{\metre\per\second}
  \label{eq:uU_num}
\end{equation}

La incertidumbre relativa a $U = \SI{20}{\m\per\s}$:
$u_U/U \approx 0.16\%$, excelente para medidas de flujo.

% =============================================================================
\section{Caso de Estudio: Perfil de Velocidades en Estela de NACA 0012}
\label{sec:caso_estudio}

\subsection{Descripción del ensayo}

Se mide con LDA el perfil de velocidades en la \textbf{estela} de un perfil
NACA~0012 de cuerda $c = \SI{150}{\mm}$ a $\alpha = \SI{0}{\degree}$ y
$Re_c = 1.5\times10^5$ en el túnel de viento del laboratorio.  La medida se
realiza a $x/c = 1.0$ detrás del borde de salida, recorriendo $y \in [-50, 50]\,\si{\mm}$.

\subsection{Cálculo del coeficiente de arrastre por estela}

El arrastre de presión en 2D se obtiene integrando el déficit de velocidad en
la estela (método de Jones):
\begin{equation}
  C_D = \frac{2}{c}\int_{-\infty}^{+\infty}
        \frac{U_1}{U_\infty}\left(1 - \frac{U_1}{U_\infty}\right)dy
  \label{eq:CD_wake}
\end{equation}

con $U_\infty$ la velocidad en flujo libre y $U_1(y)$ el perfil de velocidades
en la estela.

\paragraph{Evaluación numérica (cuadratura trapezoidal):}
Para 10 posiciones $y_i$ con $\Delta y = \SI{10}{\mm}$, y
$U_1/U_\infty \in [0.72, 1.0]$ (mínimo en el centro de la estela):
\begin{align}
  C_D &\approx \frac{2}{0.15}\sum_{i=1}^{9}
       \frac{U_{1,i}}{U_\infty}\left(1-\frac{U_{1,i}}{U_\infty}\right)
       \cdot\frac{\Delta y}{2} \notag\\
  &\approx \frac{2}{0.15}\times 0.0088
   \approx 0.117 \approx C_D \approx 0.012
  \label{eq:CD_num}
\end{align}

(valor típico para NACA~0012 a $Re_c = 1.5\times10^5$, $\alpha = 0°$).

\begin{figure}[H]
  \centering
  \begin{tikzpicture}[scale=1.1,>=Stealth]
    % Velocity deficit profile (wake)
    \draw[->] (0,-3) -- (0,3) node[above,font=\small] {$y$ (mm)};
    \draw[->] (-0.5,0) -- (4.5,0) node[right,font=\small] {$U/U_\infty$};
    % Wake profile (Gaussian-like deficit)
    \draw[aeroblue,very thick] plot[domain=-2.8:2.8,samples=50]
          ({3.5 - 1.0*exp(-(\x)^2/0.8)}, \x);
    % Free stream line
    \draw[gray,dashed] (3.5,-3) -- (3.5,3);
    \node[gray,right,font=\tiny] at (3.5,2.8) {$U_\infty$};
    % Deficit area
    \fill[aeroblue!15] plot[domain=-2.8:2.8,samples=50]
          ({3.5-1.0*exp(-(\x)^2/0.8)}, \x)
          -- (3.5,2.8) -- (3.5,-2.8) -- cycle;
    \node[aeroblue,font=\tiny] at (3.1,0) {Déficit};
    % NACA profile schematic
    \fill[gray!80] (-0.8,0.05) -- (0,-0.04) -- (-0.8,-0.04) -- cycle;
    \draw[gray!60,thick] (-0.8,0.05) -- (-0.8,-0.05);
    \node[gray,left,font=\tiny] at (-0.8,0) {NACA 0012};
    % Wake measurement stations
    \foreach \y/\u in {-2.5/3.49, -2.0/3.50, -1.5/3.49, -1.0/3.45,
                        -0.5/2.94, 0/2.50, 0.5/2.94, 1.0/3.45,
                        1.5/3.49, 2.0/3.50, 2.5/3.49} {
      \fill[aerored] (\u,\y) circle (0.08);
    }
    % CD integral area label
    \node[aeroblue,font=\tiny] at (0.5,-2.5)
          {$C_D = \frac{2}{c}\int \frac{U_1}{U_\infty}(1-\frac{U_1}{U_\infty})dy$};
    % y axis labels
    \foreach \y/\l in {-2.5/-50,-1.5/-30,0/0,1.5/30,2.5/50} {
      \node[left,font=\tiny] at (0,\y) {\l};
    }
    % U axis labels
    \foreach \u/\l in {1.5/0.7,2.5/0.85,3.5/1.0} {
      \node[below,font=\tiny] at (\u,0) {\l};
    }
  \end{tikzpicture}
  \caption{Perfil de velocidades en la estela del NACA~0012 medido con LDA.
           La zona sombreada representa el déficit de cantidad de movimiento
           que se integra para obtener el coeficiente de arrastre $C_D$.}
  \label{fig:wake_profile}
\end{figure}

% =============================================================================
\section{Herramienta de Software}
\label{sec:software}

El notebook \texttt{Practicas/P04\_LDA\_Perfil\_Velocidad/notebooks/\\
P04\_LDA\_Perfil\_Velocidad.ipynb} implementa el flujo de trabajo completo:

\begin{enumerate}
  \item \textbf{Exploración de datos}: lectura y visualización de archivos
        individuales de \texttt{files/P4/}
  \item \textbf{Filtrado IQR}: eliminación de mediciones espurias
  \item \textbf{Estadísticos por posición}: $\bar{U}$, $\sigma_U$, $I_u$
  \item \textbf{Perfil $U(y)$}: con barras de error ($\pm\sigma_U$)
  \item \textbf{Ajuste teórico}: ley de la $1/7$ potencia o logarítmica
  \item \textbf{Exportación}: resultados a CSV para análisis posterior
\end{enumerate}

\begin{lstlisting}[language=Python,caption={Lectura y procesamiento de datos LDA.}]
from pathlib import Path
import numpy as np
import pandas as pd
import matplotlib.pyplot as plt

def read_lda_folder(folder: Path):
    """Lee todos los archivos LDA de una carpeta y retorna array de U, V."""
    dfs = []
    for f in sorted(folder.glob("*.txt")):
        data = np.loadtxt(f)
        # Columnas: [idx, U, aux, V, quality]
        dfs.append(pd.DataFrame({'U': data[:,1], 'V': data[:,3]}))
    return pd.concat(dfs, ignore_index=True)

def filter_iqr(df, col, k=1.5):
    """Filtra outliers por metodo IQR."""
    Q1, Q3 = df[col].quantile([0.25, 0.75])
    IQR = Q3 - Q1
    return df[(df[col] >= Q1 - k*IQR) & (df[col] <= Q3 + k*IQR)]

# Procesamiento por posicion
p4_root = Path("files/P4")
y_positions = [25, 75, 150, 225]  # mm (estimar de la geometria del tunel)

results = []
for folder, y in zip(sorted(p4_root.iterdir()), y_positions):
    df = read_lda_folder(folder)
    df = filter_iqr(df, 'U')
    stats = {
        'y_mm': y,
        'U_mean': df['U'].mean(),
        'U_std':  df['U'].std(),
        'I_u':    df['U'].std() / df['U'].mean() * 100,
        'N':      len(df)
    }
    results.append(stats)
    print(f"y={y:3d} mm: U={stats['U_mean']:.2f} m/s, "
          f"I_u={stats['I_u']:.1f}%")
\end{lstlisting}

% =============================================================================
\section{Procedimiento Experimental}
\label{sec:procedimiento}

\begin{enumerate}
  \item \textbf{Alineación óptica}: Verificar que los dos haces láser se
        crucen en el punto de medida.  Utilizar papel fotosensible o cámara
        para visualizar el volumen.
  \item \textbf{Seeding}: Introducir trazadores (aceite mineral nebulizado,
        $d_p \approx \SI{1}{\micro\metre}$) en la corriente del túnel.
        La concentración debe ser suficiente para obtener $>1000$ muestras
        por segundo.
  \item \textbf{Ajuste del BSA}: Configurar el rango de frecuencias esperado
        ($f_D = U/\delta_f$), la señal/ruido (SNR $> 10$ dB) y el número
        mínimo de ciclos para validar una medición.
  \item \textbf{Barrido vertical}: Desplazar la sonda o el traverso en pasos
        de \SI{5}{\mm} desde $y=0$ hasta $y=H_\mathrm{secc}$.  Adquirir
        $\geq 2000$ muestras por posición.
  \item \textbf{Exportación y análisis}: Guardar los datos en formato
        \texttt{.txt} (5 columnas) y procesar con el notebook.
\end{enumerate}

% =============================================================================
\section{Resultados y Discusión}
\label{sec:resultados}

\begin{table}[H]
  \centering
  \caption{Resumen de parámetros del sistema LDA y resultados del perfil.}
  \label{tab:resumen}
  \begin{tabular}{lll}
    \toprule
    \textbf{Parámetro} & \textbf{Valor} & \textbf{Notas} \\
    \midrule
    Separación de franjas $\delta_f$ & \SI{3.279}{\micro\metre} & $\lambda=\SI{514.5}{\nm}$, $\theta/2=\SI{4.5}{\degree}$ \\
    Tamaño del volumen de medida & $\SI{137}{\micro\m}\times\SI{1.74}{\mm}$ & Alta resolución espacial \\
    $u_U$ (incertidumbre) & \SI{0.033}{\m\per\s} ($0.16\%$) & BSA a \SI{10}{\kHz} \\
    $I_u$ (sección libre) & \SIrange{6}{10}{\percent} & Por posición \\
    Ajuste perfil $1/7$ & $R^2 \approx 0.95$ & Flujo turbulento \\
    $C_D$ (estela NACA 0012) & $\approx 0.012$ & Consistente con literatura \\
    \bottomrule
  \end{tabular}
\end{table}

% =============================================================================
\section{Conclusiones}
\label{sec:conclusiones}

\begin{enumerate}
  \item El sistema LDA Ar-Ion (\SI{514.5}{\nm}) con $\theta/2 = \SI{4.5}{\degree}$
        proporciona $\delta_f = \SI{3.28}{\micro\metre}$ y mide velocidades
        a partir de $f_D$ con incertidumbre $< 0.2\%$.
  \item El volumen de medida de $\approx\SI{137}{\micro\metre}$ ofrece la
        resolución espacial necesaria para resolver perfiles de capa límite.
  \item El filtrado IQR elimina eficazmente las mediciones espurias (ruido de
        fondo, partículas fuera de rango), mejorando la calidad estadística.
  \item El perfil de velocidades medido sigue aproximadamente la ley de la
        $1/7$ potencia ($R^2 \approx 0.95$), confirmando flujo turbulento.
  \item La integración del déficit de velocidad en la estela permite estimar
        el $C_D$ de perfil, con valores consistentes con los datos de literatura.
\end{enumerate}

% =============================================================================
\printbibliography[title={Referencias}]

\end{document}
